\chapter{作为数的分数}\label{ch:g5:fractions}

\omterm{def:g4:fractions:def}{分数}不仅仅是一张饼的图画：它是一个真正的
\emph{数}，在数轴上有自己的位置——有时
恰好与某个小数位于同一位置。本章将连接这两个
世界。（\omterm{def:g4:fractions:def}{分数}运算的完整法则见
\cref{ch:g6:fractions}及以后各章。）

\section{再论数轴上的分数}

\begin{example}[在数轴上定位与读数]\label{ex:g5:fractions:place}
在一条五等分的直线上，$\frac{7}{5}$ 是 $7$ 步之外的点（每步一份，从 $0$ 起算），
越过 $1$（即 $\frac55$），位于 $1 + \frac25$ 处。每个\omterm{def:g4:fractions:def}{分数}都在直线的某处；
\omterm{def:g4:fractions:def}{分母}相同时，\omterm{def:g4:fractions:def}{分子}越大，
位置越靠右。
\end{example}

\begin{proposition}[同一点，多个名字]\label{prop:g5:fractions:equal}
把每一份都二等分（或三等分\,\dots），
点的位置并不会改变：
\[
\frac{1}{2} = \frac{2}{4} = \frac{3}{6} = \frac{5}{10},
\qquad
\frac{2}{3} = \frac{4}{6} = \frac{20}{30}.
\]
\omterm{def:g4:fractions:def}{分子}与\omterm{def:g4:fractions:def}{分母}同时乘（或除以）同一个数，
所得\omterm{def:g4:fractions:def}{分数}是同一\omterm{def:g4:fractions:def}{分数}的另一种写法。
\end{proposition}
\admitted

\begin{omfigure}
\begin{tikzpicture}[scale=0.62]
  \foreach \i in {0,1} \draw[omExo] (\i*4,0) rectangle (\i*4+4,1.0);
  \fill[omDef!40] (0.06,0.06) rectangle (3.94,0.94);
  \draw[thick] (0,0) rectangle (8,1.0);
  \node[right, font=\small] at (8.3,0.5) {$\frac12$};
  \begin{scope}[yshift=-1.7cm]
  \foreach \i in {0,...,3} \draw[omExo] (\i*2,0) rectangle (\i*2+2,1.0);
  \fill[omDef!40] (0.06,0.06) rectangle (1.94,0.94);
  \fill[omDef!40] (2.06,0.06) rectangle (3.94,0.94);
  \draw[thick] (0,0) rectangle (8,1.0);
  \node[right, font=\small] at (8.3,0.5) {$\frac24$};
  \end{scope}
  \begin{scope}[yshift=-3.4cm]
  \foreach \i in {0,...,9} \draw[omExo] (\i*0.8,0) rectangle
    (\i*0.8+0.8,1.0);
  \foreach \i in {0,...,4} \fill[omDef!40] (\i*0.8+0.05,0.06)
    rectangle (\i*0.8+0.75,0.94);
  \draw[thick] (0,0) rectangle (8,1.0);
  \node[right, font=\small] at (8.3,0.5) {$\frac{5}{10}$};
  \end{scope}
\end{tikzpicture}

{\small 三个长条，三种名字，同一个量：$\frac12 = \frac24 =
\frac{5}{10}$.}
\end{omfigure}

\section{分数与小数}

\begin{proposition}[可以写成小数的分数]\label{prop:g5:fractions:decimal}
十分之几、百分之几、千分之几的\omterm{def:g4:fractions:def}{分数}\emph{都是}
小数：
\[
\frac{3}{10} = 0.3, \qquad \frac{27}{100} = 0.27 .
\]
另有一些\omterm{def:g4:fractions:def}{分数}也有小数写法，只需化为
十分之几或百分之几：
\[
\frac{1}{2} = \frac{5}{10} = 0.5, \qquad
\frac{1}{4} = \frac{25}{100} = 0.25, \qquad
\frac{3}{4} = 0.75, \qquad
\frac{1}{5} = \frac{2}{10} = 0.2 .
\]
但并非所有\omterm{def:g4:fractions:def}{分数}都如此：$\frac13 = 0.333\dots$ 永远不会结束——
\omterm{def:g4:fractions:def}{分数}本身才是它唯一的精确写法。
\end{proposition}
\admitted

\begin{omfigure}
\begin{tikzpicture}[scale=1.55]
  \draw[-Stealth] (-0.2,0) -- (7.6,0);
  \foreach \i in {0,...,28} \draw ({\i*0.25},-0.05) -- ({\i*0.25},0.05);
  \foreach \i in {0,4,...,28} \draw ({\i*0.25},-0.09) -- ({\i*0.25},0.09);
  \foreach \x in {0,...,7} \node[below, font=\small] at (\x,-0.1) {$\x$};
  \fill[omDef] (0.5,0) circle (1.6pt)
    node[above, font=\small] {$\frac12 = 0.5$};
  \fill[omThm] (2.75,0) circle (1.6pt)
    node[above, font=\small] {$\frac{11}{4} = 2.75$};
  \fill[omMeth] (6.25,0) circle (1.6pt)
    node[above, font=\small] {$\frac{25}{4} = 6.25$};
\end{tikzpicture}

{\small 一条数轴，两种写法：每个点都可以带一个\omterm{def:g4:fractions:def}{分数}名
和一个小数名。}
\end{omfigure}

\begin{example}[选择更方便的名字]\label{ex:g5:fractions:choose}
求 $12$ 的四分之三：用\omterm{def:g4:fractions:def}{分数}名更方便，
$\frac34$ 的 $12$ 是 $9$。比较 $\frac34$ 与 $0.8$ 时，
用小数名更方便，$0.75 < 0.8$。两种写法都会用，大有益处。
\end{example}

\section{比较分数的大小}

\begin{method}[三种比较方法]\label{met:g5:fractions:compare}
\begin{enumerate}
  \item \emph{\omterm{def:g4:fractions:def}{分母}相同}：份数多的更大，
        $\frac57 > \frac37$；
  \item \emph{\omterm{def:g4:fractions:def}{分子}相同}：\emph{\omterm{def:g4:fractions:def}{分母}越小}的\omterm{def:g4:fractions:def}{分数}更大：
        $\frac35 > \frac38$（五分之一
        比八分之一大）；
  \item \emph{参照点}：分别与 $\frac12$ 或 $1$ 比较：
        $\frac38 < \frac12 < \frac59$，因此 $\frac38 < \frac59$。
\end{enumerate}
\end{method}

\begin{example}\label{ex:g5:fractions:compare}
谁吃的比萨更多：阿里吃了 $\frac58$，还是贝亚吃了 $\frac54$？贝亚的
\omterm{def:g4:fractions:def}{分数}大于 $1$（五个\omterm{def:g3:division:half}{四分之一}比一整张比萨还多），
阿里的\omterm{def:g4:fractions:def}{分数}小于 $1$：贝亚吃得更多——事实上她
吃了一整张比萨还多。
\end{example}

\section{练习}

\begin{exercise}[$\star$]\label{exo:g5:fractions:1}
画一条从 $0$ 到 $2$ 的数轴，每个单位五等分，并标出：
$\frac25$; \ $\frac55$; \ $\frac85$; \ $\frac{10}{5}$.
\end{exercise}

\begin{exercise}[$\star$]\label{exo:g5:fractions:2}
补全下列等值\omterm{def:g4:fractions:def}{分数}：
$\frac12 = \frac{?}{8}$; \quad
$\frac34 = \frac{?}{8}$; \quad
$\frac{2}{5} = \frac{4}{?}$; \quad
$\frac{6}{9} = \frac{2}{?}$.
\end{exercise}

\begin{exercise}[$\star$]\label{exo:g5:fractions:3}
哪些\omterm{def:g4:fractions:def}{分数}可以化为整数？如果可以，写出该整数：
$\frac82$; \quad $\frac{9}{4}$; \quad $\frac{15}{3}$; \quad
$\frac{20}{5}$.
\end{exercise}

\begin{exercise}[$\star$]\label{exo:g5:fractions:4}
化为小数：$\frac{7}{10}$；\quad $\frac{41}{100}$；\quad
$\frac12$; \quad $\frac34$; \quad $\frac15$.
\end{exercise}

\begin{exercise}[$\star$]\label{exo:g5:fractions:5}
化为\omterm{def:g4:fractions:def}{分数}（十分之几或百分之几）：$0.9$；\quad $0.13$；
\quad $2.5$；\quad $0.75$（再说出它是几个\omterm{def:g3:division:half}{四分之一}）。
\end{exercise}

\begin{exercise}[$\star$]\label{exo:g5:fractions:6}
比较下列各组\omterm{def:g4:fractions:def}{分数}，并指出所用的方法
（见\cref{met:g5:fractions:compare}）：
\[
\frac47 \;?\; \frac67, \qquad
\frac35 \;?\; \frac37, \qquad
\frac25 \;?\; \frac{7}{12}, \qquad
\frac98 \;?\; \frac{11}{12} .
\]
\end{exercise}

\begin{exercise}[$\star$]\label{exo:g5:fractions:7}
借助小数名称，将下列各数排序：$\frac12$；\ $0.4$；\ $\frac34$；\
$0.6$; \ $\frac15$.
\end{exercise}

\begin{exercise}[$\star$]\label{exo:g5:fractions:8}
某班有 $28$ 名学生，其中\omterm{def:g3:division:half}{四分之一}的学生学乐器，
\omterm{def:g3:division:half}{一半}的学生参加体育运动。两组各有多少名学生？哪一组人数更多？
\end{exercise}

\begin{exercise}[$\star$]\label{exo:g5:fractions:9}
将每个\omterm{def:g4:fractions:def}{分数}与一个时间配对：$\frac14$ 小时，$\frac12$ 小时，
$\frac34$ 小时，$\frac{1}{60}$ 小时——对应 $30$ 分钟，$45$ 分钟，$1$ 分钟，
$15$ 分钟。
\end{exercise}

\begin{exercise}[$\star\star$]\label{exo:g5:fractions:10}
汤姆喝掉了 $\frac35$ 的饮料，他的瓶子共 $50$ 厘升；安娜喝掉了 $\frac12$ 的饮料，
她的瓶子共 $60$ 厘升。谁喝得更多？（请分别以厘升为单位算出两人的饮用量——
仅比较\omterm{def:g4:fractions:def}{分数}在这里是不够的，请说明理由。）
\end{exercise}

\begin{exercise}[$\star\star$]\label{exo:g5:fractions:11}
写出一个严格介于 $\frac12$ 与 $1$ 之间的\omterm{def:g4:fractions:def}{分数}；再写出一个
严格介于 $\frac12$ 与 $\frac34$ 之间的\omterm{def:g4:fractions:def}{分数}。（借助小数写法
可能会有帮助。）
\end{exercise}
